\section{Conclusion}

We formulated many-to-one DRT service-menu design as an assortment optimization problem over meeting-point and pickup-window bundles. The implemented framework combines MNL passenger choice, predicted route-aware bundle costs, common-offset pricing, exact enumeration for small candidate sets, and greedy forward selection for larger candidate sets. It also logs exact-vs-greedy approximation diagnostics so that the scalable solver can be evaluated against a surrogate benchmark.

The empirical contribution is a disciplined evidence program rather than a universal dominance claim. Exact-vs-greedy diagnostics quantify the cost of scalable menu construction; robust ETA-filtering comparisons show how eligibility rules reshape the feasible assortment before optimization acts; and low-, medium-, and high-uptake regimes clarify when profit, acceptance, non-home uptake, and welfare comparisons are behaviorally meaningful. Austin and Seattle remain external directional checks under a low-sample constraint.

The main practical lesson is that service-menu optimization should be managed as a layered system. In the studied settings, ETA filtering was the most impactful layer; generalization to other demand regimes and network structures is an open question. Operators considering service-menu optimization should audit filtering and approximation quality before interpreting pricing or welfare comparisons, but the current evidence supports these recommendations only within the tested calibration range. Broader validation with additional city splits, empirically estimated demand parameters, and richer real-world-derived instances remains the next step.
